%% Šablona pro psaní závěrečných prací FBMI ČVUT
% Upraveno na základě požadavků na závěrečné práce schválených dne 24.2.2020.
%
% Vytvořeno ve spolupráci členů BRAIN Teamu.
% 
% Šablonu pro vás upravil/a:    Ing. Václava Piorecká, Ph.D.
% Kontakt:                      vaclava.piorecka@fbmi.cvut.cz

\documentclass[a4paper,12pt]{article}   % Definice - velikost dokumentu, základní velikost písma, typ
\usepackage[a4paper, top=2.5cm, left=3.5cm, right=2.5cm, bottom=2.5cm]{geometry}		% nastavení okrajů

%% Přidání balíčků podporujících různé funkcionality - ZÁKLADNÍ
\usepackage{amsmath,float}				% balíček pro matiku
\usepackage[dvipsnames]{xcolor}
\usepackage{color}
\usepackage{footnote}					% poznámka pod čarou
\usepackage{url}						% url odkazy
\DeclareUrlCommand\url{\def\UrlLeft{<}\def\UrlRight{>} \urlstyle{same}}   % nastavení URL odkazů pro podle normy ČSN ISO 690

\usepackage{float}						% plovoucí prostředí
\usepackage[utf8]{inputenc}			    % kódování	
\usepackage[czech]{babel}				% čeština
\usepackage{enumerate}      			% seznamy 
\usepackage{amsfonts}      				% množiny Z,R,N dvojitě
\usepackage{amssymb}      				% znaky úhlu a tak
\usepackage[pdftex]{graphicx}
\usepackage{setspace}
\usepackage{multicol}					% tabulka, slučování sloupců
\usepackage{multirow}                   % tabulka, slučování řádků
\usepackage{fancyhdr}					% záhlaví a zápatí stránky
\usepackage{chngcntr}                   % číslování (číslování rovnic, obrázků dle kapitol)
\usepackage{array}                      % rozšíření práce s tabulkami
\usepackage{helvet}                     % předefinuje \sfdefault to uhv (pro úvodní stránku)
\usepackage[flushleft]{threeparttable}  % prostředí pro tabulky - přidává vysvětlující poznámky pod tabulku

% Následující balíčky řeší citace v dokumentu. 
\usepackage{csquotes}                       
\usepackage[style=iso-numeric]{biblatex}     
\addbibresource{citace.bib}             % zdrojový soubor s citacemi

%% Přidání balíčků podporujících různé funkcionality - VOLITELNÉ, DOPLŇKOVÉ
% Existuje celá řada dalších
%\usepackage[]{algorithm2e}				    % balíček pro pseudokód
%\usepackage{ifthen}                        % pro algoritmy if else
%\usepackage{paralist}                      % Rozšířená možnost pro seznamy. Větší škála a možnosti jak seznamy dělat autoamticky. 
%\usepackage{fontspec}                      % specifické fonty
%\usepackage{icomma}      				    % není mezera za desetinnou čárkou
% \usepackage[titletoc,title]{appendix}     % automatické vytvoření příloh

%% Přidání speciálních příkazů 
\newcommand{\at}{\makeatletter @\makeatother}           % Vytiskne zavináč - \at
\newcommand{\degree}[1][]{\ensuremath{{#1}^\circ}}      % Vytiskne stupeň - \degree
\newcolumntype{C}[1]{>{\centering\let\newline\\\arraybackslash\hspace{0pt}}m{#1}}                                                     % zarovnání v tabulce, vycentrování, potřebuje \usepackage{array}


%% Ostatní definice a nastavení
\clubpenalty 10000		% penalizace sirotků, sirotek: poslední řádek odstavce je na nové stránce
\widowpenalty 10000		% penalizace vdov, vdova: první řádek nového odstavce je na konci stránky



\DeclareGraphicsExtensions{.pdf,.png,.jpg}	    % nahrávání obrázků, rošíření
\graphicspath{{obrazky/}} 				        % umístění obrázků

\counterwithin{figure}{section}         % číslování obrázků dle sekcí/kapitol
\counterwithin{table}{section}          % číslování tabulek dle sekcí/kapitol
\numberwithin{equation}{section}        % číslování rovnic dle sekcí/kapitol


\setlength{\parskip}{6pt}                   % odsazení mezi odstavci
\setlength{\parindent}{0.75cm}              % odsazení odstavců od okraje
\renewcommand{\baselinestretch}{1.20}	    % řádkování 1,2 - odpovídá pevnému řádkování 17 bodů

\usepackage{tocloft}                    % Nastaví tučné zvýraznění sekcí v obsahu
\renewcommand{\cftsecleader}{\cftdotfill{\cftdotsep}}   % přidá vodící linku do obsahu u sekcí


%% Ostatní neimplementované, pouze návod jak případně doplnit
% Vytvořit seznam použitých algoritmů a přejmenování názvu objektu na Algoritmus.
% \usepackage[algosection]{algorithm2e}
% \SetAlgorithmName{Algoritmus}{algorithm}{Seznam algoritmů}
   

%% Deklarace názvů - přepsat dle autora
\newcommand{\autor}{Monika Pulcová}
\newcommand{\vedouci}{Jméno Příjmení - vedoucího práce včetně správných titulů}
\newcommand{\nazevENG}{Název v jazyce práce}
\newcommand{\nazev}{Automatické určení tvaru a velikosti nanočástic pro biomedicínské aplikace}
\newcommand{\typ}{Semestrální projekt 2}
\newcommand{\rok}{2022/23}
% Pro program a obor jsou instrukce zde: http://www.fbmi.cvut.cz/fakulta/uredni-deska 
% Nově akreditované mají pouze studijní programy
\newcommand{\program}{Biomedicínská technika}

%% Vlastní začátek dokumentu
\begin{document}
    
	\pagestyle{empty}

	\begin{titlepage}
 		\begin{center}
 		\begin{figure}[!h]
			\centering
 			\includegraphics[width=0.2\textwidth]{symbol_cvut_konturova_verze}
 		\end{figure}
 		\textsf{\large{\textbf{ČESKÉ VYSOKÉ UČENÍ TECHNICKÉ V PRAZE}}}\\
 		%\vspace{0pt}   		
         {\color{NavyBlue}\makebox[\linewidth]{\rule{\textwidth}{0.4mm}}}
         %{\color{NavyBlue}\hrule }  	    \vspace{6pt}
 		\textsf{\normalsize{\textbf{FAKULTA BIOMEDICÍNSKÉHO INŽENÝRSVÍ}}}\\
		\textsf{\textbf{Katedra biomedicínské techniky}}\\	
		
		\vfill
 		
		\textsf{\Large{\textbf{\nazev}}}\\
	    \vspace{24pt}
		\textsf{\Large{\textbf{\nazevENG}}}\\
		\vspace{24pt}
		\textsf{\typ}\\ 
		\vfill
		\end{center}
		\textsf{Studijní program: \program} \\
		\textsf{Vedoucí práce: \vedouci}\\
				
		\begin{center}
		\textsf{\textbf{\autor}} \\ [0.5cm] 
		
		{\color{NavyBlue}\makebox[\linewidth]{\rule{\textwidth}{0.4mm}}} 
 		
		\textsf{\textbf{Kladno \rok}}
		\end{center}
			

		\clearpage


	\end{titlepage}

    
	\null\vfill
	\section*{Prohlášení}
	% \hspace ruší odsazení odstavce - v šabloně u prohlášení odstavce odsazené nejsou. 
    Prohlašuji, že jsem bakalářskou práci s názvem \uv{\nazev} vypracoval/a samostatně a~použil/a k~tomu úplný výčet citací použitých pramenů, které uvádím v seznamu přiloženém k~bakalářské práci.

    \hspace{-0.75cm}Nemám závažný důvod proti užití tohoto školního díla ve smyslu \S 60 Zákona č.121/2000~Sb., o~právu autorském, o~právech souvisejících s právem autorským a~o~změně některých zákonů (autorský zákon), ve znění pozdějších předpisů. 
    
    \vspace{1em}
    
    \hspace{-0.75cm}V Kladně dne \ldots \ldots \ldots \hfill \ldots \ldots \ldots \ldots \ldots \ldots \ldots \ldots \ldots \ldots

    \hspace{10cm} \textbf{\autor}

	\clearpage
	
	\null\vfill
	\section*{Poděkování}
	Rád/a bych poděkoval/a… 
	    
    Poděkování je nepovinné, ale obvyklé. 
    Vedoucímu práce se zpravidla děkuje, oponentovi zásadně ne. 
    Poraďte se s vedoucím práce, zda by nebylo vhodné uvést v poděkování číslo grantu, ze kterého byla práce podpořena.  

    \clearpage	

	
	\null\vfill	
	\section*{ABSTRAKT}
        \subsection*{\nazev:}
	    Výstižná charakteristika cílů, metod, výsledků, (diskuse) a závěrů bakalářské práce v rozsahu asi 10 řádků. 
	\subsection*{Klíčová slova}
		výčet tří až pěti klíčových slov nebo sousloví charakterizujících obsah bakalářské práce
	\clearpage
		
		
	\null\vfill	
	\section*{ABSTRACT}
        \subsection*{\nazevENG:}
		 
        A concise summary of aims, methods, results, discussion (if needed) and conclusions of the Bachelor’s Thesis within the range of about 10 lines.
    
	\subsection*{Key words}
		Listing 3 to 5 key words characterizing the subject-matter of the Bachelor's Thesis.
	\clearpage
	
    \pagestyle{plain}	% Číslování stránek začíná odsud
	
	\tableofcontents			% Vloží obsah

	\clearpage

	\section*{Seznam symbolů a zkratek} %sekce = nadpis, s * neni v obsahu
	\addcontentsline{toc}{section}{Seznam symbolů zkratek}
	\subsection*{Seznam symbolů}

\begin{table}[h]
	\label{tab:symboly}
	\catcode`\-=12          % Tento řádek je tam kvůli použití cline pro czech babel. Jinak to bere pomlčku jako znak a nevnímá ji jako rozsah. 
	\begin{center}
		\begin{tabular}{p{2.5cm}p{2.5cm}p{9.25cm}}
			\noalign{\hrule height 2pt}
			Symbol  & Jednotka              & Význam \\
			\noalign{\hrule height 2pt}
			$CDP_{end}$ & cmH$_2$O          & Střední distenzní tlak po ustálení přechodového děje \\     
			$f_{HFO}$   & Hz                & Frekvence vysokofrekvenčních oscilací \\
			$N$         & 1$\backslash$s    & Parametr filtru typu dolní propust připojeného k~PID regulátoru\\
			\noalign{\hrule height 2pt}
	    \end{tabular}
	\end{center}
\end{table}

\subsection*{Seznam zkratek}
\begin{table}[h]
	\label{tab:zkratky}
	\catcode`\-=12          % Tento řádek je tam kvůli použití cline pro czech babel. Jinak to bere pomlčku jako znak a nevnímá ji jako rozsah. 
	\begin{center}
		\begin{tabular}{p{2.5cm}p{12.25cm}}
			\noalign{\hrule height 2pt}
			Zkratka  & Význam              \\
			\noalign{\hrule height 2pt}
			EEG     & Elektroencefalograf                       \\
			ALI	    & Akutní plicní selhání (Acute Lung Injury) \\
            PID	    & Proporcionálně-integračně-derivační       \\
			\noalign{\hrule height 2pt}
	    \end{tabular}
	\end{center}
\end{table}
	\clearpage
		
	    %\addcontentsline{toc}{section}{Seznam tabulek}
		%\listoftables 		% seznam tabulek
	
		%\clearpage 			% kones stránky a odskok na další
		
		%\addcontentsline{toc}{section}{Seznam obrázků}
		%\listoffigures 		% seznam obrázků
		%\clearpage
	
		%\addcontentsline{toc}{section}{Seznam algoritmů}
		%\listofalgorithms
		%\clearpage
	
	
	\section{Úvod}
	Úvod obsahuje nejprve stručný obecný úvod do řešené problematiky (definuje oblast, kterou se práce zabývá, uvádí motivaci apod.). 
Obecný úvod má svým rozsahem tvořit velmi malou část celé práce.
	\clearpage
	
	\section{Přehled současného stavu}
	Přehled aktuálního stavu řešené problematiky podrobně shrnuje (1) současný stav poznání a výchozí podmínky pro řešení a to jak v tuzemsku, tak i v zahraničí (2) definuje problém, který je nutno a který se bude v práci řešit. 
Tato část práce je převážně vytvořena jako rešerše za použití mnoha literárních zdrojů. 
Při výkladu se postupuje od obecnějších informací k informacím co nejkonkrétnějším a od toho, co se o dané problematice ví, k tomu, co je neznámé a aktuálně vhodné k řešení.
Z~takto uspořádaného výkladu pak logicky vyplynou cíle práce vytyčené níže.


	\clearpage
	
    \section{Cíle práce}
	Zde mají být výstižně popsány vytyčené cíle Vaší práce, vycházející ze zadání práce. 
Na rozdíl od velmi stručného zadání práce je nutné cíle v této části specifikovat podrobněji a vhodné je i rozvést cíle do specifických podcílů.
	\clearpage
	
	\section{Metody}
	Kapitola obsahuje detailní popis způsobu řešení problému studentem. 
V závislosti na charakteru řešeného problému je tuto část práce možné rozdělit do více kapitol/podkapitol, kdy názvy kapitol jsou voleny konkrétněji s ohledem na jejich obsah.

Popisovány jsou postupy aplikované k dosažení výsledků práce a rovněž např. použité přístroje a materiál, metody zpracování dat a jejich statistického vyhodnocení apod. 
V případě měření s živými subjekty tato část práce obsahuje informaci, jak byly ošetřeny etické otázky výzkumu a charakteristiku subjektů dle zvyklostí v biomedicínských časopisech.
Struktura a obsah této části je detailně probírána a~procvičována v příslušných seminářích na oboru BMT.
V případě, že text obsahuje matematický vzorec, na který se bude text později odkazovat, uvádějte vzorec na samostatném řádku, vycentrovaný na střed řádku a s číslem, které udává pořadí mezi číslovanými vzorci v kapitole, jako je tomu v příkladu vztahu pro elektrický odpor $R$:
\begin{equation}
	\label{rce:odpor}
	R = \frac{U}{I}
\end{equation}
kde $U$ je napětí a $I$ je proud. 
Pokud je vzorec součástí věty, jako v předchozím vztahu~\ref{rce:odpor}, pokračujte za ním textem bez odsazení nového odstavce. 
Vzorce vkládejte pomocí možností editoru, nekopírujte vzorce z jiných pramenů.  

\subsection{Ukázka citací}
Ukázka citací:
Byl využit software Matlab \cite{MATLAB} s doplňkovým toolboxem EEGLAB~\cite{eeglab}. 
Princip algoritmu k-means najdete například v publikaci~\cite{Krajca2011}.

Ukázka více citací~\cite{MATLAB, eeglab}.
	\clearpage
	
	\section{Výsledky}
	Věnujte kapitolu pouze přehlednému podání výsledků, nikoliv jejich diskusi. 
Data uvádějte zejména v grafech a tabulkách. 
Preferovány jsou grafy – tabulky se všemi naměřenými hodnotami, ze kterých grafy vycházejí, lze umístit do příloh práce.

Výsledky mají vždy obsahovat hlavní text, který zasadí prezentované obrázky a tabulky do souvislosti s předchozím textem a čtenáře prezentovanými daty provede. 
Prezentování výsledků ve formě nekomentovaného obrázkového alba je v drtivé většině případů nevhodné.
Struktura a obsah této části je detailně probírána a procvičována v příslušných seminářích na oborech BMT, BME a SIPZ.

Na každý obrázek musí být uveden odkaz v textu, který má formát jako v~ná\-sle\-du\-jí\-cí větě. 
Obrázek se vždy čísluje a popisuje pod obrázkem, viz příklad na Obrázku~\ref{img:tulipany}.

\begin{figure}[h]
	\label{img:tulipany}
	\begin{center}
		\includegraphics[width=0.3\textwidth]{tulipany}
		\caption{Tulipány před ozářením kryptonitem. Fotografie: autor.}
	\end{center}
\end{figure}
 
 
Obrázky číslujte podle hlavní kapitoly, ve které se vyskytují. Podkapitoly se již neuvažují. 
To znamená, že obrázky v úvodu (typicky kap. 1) budou: Obr.~1.1, Obr.~1.2 atd. V metodách (typicky druhá velká kapitola) budou číslovány Obr.~2.1, Obr.~2.2, Obr.~2.3 atd. 
Podrobněji viz dokument Často kladené dotazy týkající se psaní diplomové práce, dostupný na:
\url{https://predmety.fbmi.cvut.cz/cs/17pmbds2}

Popis tabulky, na rozdíl od obrázku, je zpravidla nad tabulkou, viz Tabulka \ref{tab:tabulka}. 
Není nutné v něm opisovat celý obsah záhlaví tabulky, které následuje hned vzápětí. 
Jednotlivé proměnné v tabulce jsou řazeny do sloupců. 
V tabulce jsou nezávislé proměnné, kategorie probandů apod. řazeny vlevo, závislé proměnné vpravo. 
Jednotky uvádějte v kulatých závorkách v záhlaví tabulky, ne u každého čísla zvlášť. 
Vysvětlující poznámky (např. dosažená hladina významnosti, zda jsou data udávána jako průměr a směrodatná chyba průměru, jaký statistický test byl použit apod.) jsou umisťovány pod tabulku a odkaz na ně se udává jako horní index (symboly, čísla, písmena) na příslušném místě tabulky. 
Na každý obrázek a tabulku je třeba odkazovat z hlavního textu.

\begin{table}[h]
	\label{tab:tabulka}
	\catcode`\-=12          % Tento řádek je tam kvůli použití cline pro czech babel. Jinak to bere pomlčku jako znak a nevnímá ji jako rozsah. 
	\begin{center}
	\begin{threeparttable}      % Pokud nepotřebujete vysvětlení pod tabulkou, stačí použít tabular (vypustit threeparttable a tablenotes).
		\caption[Tady může být kratší popisek pro seznam tabulek.]{Reakční čas $T_{20}$ signálu periferní saturace kyslíkem, $Sp$O$_2$, měřený třemi různými přístroji.}
		\vspace{1ex}
		\begin{tabular}{lccc}
			\noalign{\hrule height 2pt}
			                & \multicolumn{3}{c}{$T_{20}(s)$}                           \\  \cline{2-4}
			Fáze            & Root Radical-7    & Nellcor N-600     & Carescape B650    \\ \hline
			Hypoxická	    & $52\pm15^{*}$	    & $65\pm19^{*}$     & $56\pm15$         \\
            Hyperoxická	    & $43\pm14$	        & $55\pm28$	        & $49\pm15$         \\
            Hyperkapnická	& $75\pm23$	        & $119\pm47^{\sharp}$	& $73\pm41^{\sharp}$  \\   \noalign{\hrule height 2pt}
	    \end{tabular}
	    \begin{tablenotes}
            %\small
            \item Data byla měřena pro shodnou skupinu 14 probandů a jsou uvedena jako aritmetický průměr $\pm$ směrodatná odchylka. Symboly $^{*}$ a $^{\sharp}$ značí statisticky významný rozdíl $ \left( p \leq 0,05 \right) $ časů pro shodnou fázi.
            \end{tablenotes}
	\end{threeparttable}
	\end{center}
\end{table}



	\clearpage
	
	\section{Diskuse}
	V této části shrňte získané výsledky (hlavní zjištění práce) a následně tyto výsledky interpretujte s ohledem na cíle práce. 
Lze též získané výsledky a výstupy konfrontovat s výsledky a výstupy jiných autorů, výrobky jiných společností apod. 
Nezbytné je správné uvádění zdrojů (citace prací, které jsou zde porovnávány a diskutovány). 
Diskutují se rovněž limitace práce. 
Nakonec lze nastínit další směřování práce do budoucna, opatrně spekulovat o klinickém významu práce apod.

Struktura a obsah této části je detailně probírána a procvičována v příslušných seminářích na oboru BMT. 

	\clearpage
	
	\section{Závěr}
	Závěr stručně shrnuje splnění vytyčených cílů práce. 
Struktura a obsah této části je detailně probírána a procvičována v příslušných seminářích na oboru BMT.

Shrnutí splněných (nebo snad v menší míře nesplněných) cílů práce navazuje na úvod práce. 
Z dalších částí práce (metody, výsledky a diskuse) je přebíráno jen to nejpodstatnější a v minimální nutné míře – závěr práce nemá být zopakováním abstraktu, výsledků nebo diskuse. 
Konec závěru může obsahovat podložené spekulace o významu práce do budoucna nebo výrazná doporučení pro další výzkum nebo praxi, pokud z výsledků předkládané práce přímo vyplývají.

	\clearpage

%-------------Literatura-------------------
    \clearpage	
    \renewcommand{\refname}{Seznam použité literatury} 	% Přejmenování Reference
    \addcontentsline{toc}{section}{Seznam použité literatury}     % Přidání této kapitoly do obsahu
    \printbibliography
    \clearpage

%-------------Přílohy----------------------
    \section*{Příloha A: Požadavky na formátování práce}
    \label{app:pozadavky}
    \addcontentsline{toc}{section}{Příloha A: Požadavky na formátování práce}     % Přidání této kapitoly do obsahu

    \begin{itemize}
        \item Pro hlavní text práce používejte patkové písmo (Times New Roman, Georgia, Garamond apod.), velikost 12. Rovnice, matematické symboly apod. by měly být sázeny stejným, nebo alespoň co nejpodobnějším písmem stejné velikosti. Popisy obrázků a tabulek sázejte stejným písmem se zmenšenou velikostí. Nadpisy, čísla stránek, případné záhlaví či zápatí apod. mohou být sázeny buď stejným písmem jako hlavní text, nebo písmem bezpatkovým (např. Calibri). \textbf{V celé práci musí být použity maximálně dvě různá písma.}
        \item Řádkování práce, odsazení odstavců, velikosti písma v nadpisech apod. definují přímo jednotlivé styly použité v této šabloně.
        \item Okraje stránek práce jsou vždy 2,5 cm na každé straně plus 1 cm u hřbetu práce (levá strana).
        \item Práce je tištěna jednostranně, na papír formátu A4.
        \item Stránky se číslují arabskými číslicemi počínaje první (titulní) stranou. Číslování stránek se zobrazuje až od první stránky obsahu, což znamená, že na titulní straně, v zadání, prohlášení, poděkování a abstraktech se číslo stránky neuvádí.
        \item Hlavní kapitoly práce, počítaje v to Úvod a Závěr, jsou číslovány arabskými číslicemi. Seznam použité literatury číslo nemá. Přílohy označujte velkými písmeny anglické abecedy.
        \item Každou hlavní kapitolu práce (nadpis 1. úrovně) začínejte na samostatné stránce.
    \end{itemize}
\clearpage

\section*{Příloha B: Základní typografické zásady}
    \label{app:typo}
    \addcontentsline{toc}{section}{Příloha B: Základní typografické zásady}     % Přidání této kapitoly do obsahu
    
    \begin{itemize}
        \item Fyzikální a fyziologické veličiny a matematické proměnné se sázejí proloženě (kurzívou). Zkratky a symboly, pod kterými se neskrývá číselná hodnota, jsou sázeny normálním písmem – stejně jako označení fyzikálních jednotek.
        \item Jednotky veličin a symboly (například procenta) se v textu od číselných údajů oddělují nezlomitelnou mezerou. Zápis bez mezery má význam přídavného jména. Např. $10\,\Omega$ čteme \uv{deset ohmů} a $10\Omega$ čteme \uv{desetiohmový}.
        \item Nezlomitelnou mezeru je nutné v editoru textu vyznačit. Např. v aplikaci Microsoft Word se použije kombinace Shift, Ctrl, mezerník. 
        \item Neslabičné předložky a spojky (netýká se \uv{a}) nesmí zůstat na konci řádku. Proto za nimi používejte nezlomitelnou mezeru.
        \item Rozlišujte spojovník a pomlčku. Spojovník je krátká čára používaná ke spojení dvou slov (např. česko-anglický slovník). Pomlčka slouží k vyznačení prodlevy v textu, pak ji obvykle píšeme s mezerami, nebo k vyznačení rozsahu (5–10), kdy se píše bez mezer.
        \item Pro podrobnější informace k typografii doporučujeme např. dokument prof. Roubíka Fyzikální veličiny a číselné údaje, dostupný na stránce:
\url{https://predmety.fbmi.cvut.cz/cs/17PMBPIZ}
a dokument Jany Borůvkové Jak napsat bakalářskou práci, dostupný na stránce:
\url{http://is.mendelu.cz/dok_server/slozka.pl?id=53294;download=160152;lang=cz}
        \item Pro zajištění jazykové správnosti práce doporučujeme konzultovat Internetovou jazykovou příručku Ústavu pro jazyk český Akademie věd ČR dostupnou z: 
\url{http://prirucka.ujc.cas.cz/}
    \end{itemize}
    
\clearpage

\section*{Příloha C: Další doporučení pro přehlednost textu}
    \label{app:doporuceni}
    \addcontentsline{toc}{section}{Příloha C: Další doporučení pro přehlednost textu}     % Přidání této kapitoly do obsahu
    
    \begin{itemize}
        \item Obrázky a tabulky sázejte v textu samostatně, bez obtékání textu po stranách. Nevkládejte obrázky a tabulky na stránku před skončením odstavce. Zkontrolujte, že popis obrázku nebo tabulky zůstal na stejné straně jako vlastní obrázek nebo tabulka.
        \item První řádek odstavce by neměl zůstat sám na konci řádky (tzv. vdova) a poslední řádek odstavce by neměl zůstat sám na začátku nové stránky (tzv. sirotek).
        \item Veškeré zkratky, s výjimkou těch nejznámějších jako DNA, by měly být v práci vysvětleny při prvním výskytu v hlavním textu a současně také v abstraktu, pokud je nutné je v něm použít.
        \item Na rovnice odkazujte jejich číslem, a to až za jejich uvedením v textu práce.
        \item Všechny obrázky a tabulky v práci musí být odkazovány z hlavního textu pomocí svých čísel. 
    \end{itemize}

\clearpage

\section*{Příloha D: Obsah přiloženého CD/DVD}
    \label{app:obsah}
    \addcontentsline{toc}{section}{Příloha D: Obsah přiloženého CD/DVD}     % Přidání této kapitoly do obsahu
    
    Poslední přílohou práce je obsah přiloženého datového nosiče. 
    Typ a povinný obsah datového nosiče je specifikován na stránkách FBMI ČVUT v Praze (\url{https://www.fbmi.cvut.cz/cs/student/bakalarske-diplomove-prace}).
    
    Dále na datový nosič umístěte přílohy, které není možné pro jejich rozsah nebo charakter umístit do výtisku práce, ale které mohou být důležité pro posouzení úplnosti a kvality splnění zadání práce, jako jsou různé konstrukční výkresy, zdrojový kód programů pro zpracování naměřených dat apod.
    
\end{document}